\chapter{Problemformulierung}

\begin{tcolorbox}[colback=blue!5!white, colframe=blue!80!black, title=Kontext]
    \begin{table}[H]
        \centering
        \renewcommand{\arraystretch}{1.2}
        \begin{tabular}{@{}r|p{0.65\linewidth}@{}}
            \textbf{Gymnasialstufe} & 2. Jahr Kurzzeitgymnasium (1. Semester)                                                                                   \\
            \midrule
            \textbf{Alter}          & 16–17 Jahre                                                                                                               \\
            \midrule
            \textbf{Vorwissen}      & Grundlegende Programmierkenntnisse in Python; Einführung in Netzwerke, IP-Adressen, Ports und Client-Server-Architekturen \\
        \end{tabular}
    \end{table}
\end{tcolorbox}

Die Lernenden begegnen täglich Formen digitaler Kommunikation – sei es über WhatsApp, Instagram oder Online-Games. Doch wie Nachrichten in Sekundenbruchteilen von einem Gerät zum anderen gelangen, bleibt für sie ein „Black Box“-Phänomen. Das Problem, an dem sie sich abarbeiten sollen, lautet daher:

\begin{quote}
    \textit{„Wie ist es möglich, dass ich in Echtzeit mit anderen über das Internet kommunizieren kann?“}
\end{quote}

\textbf{Gegenwartsbedeutung:} Jugendliche nutzen Chats, Videocalls oder Multiplayer-Spiele selbstverständlich, ohne die dahinterliegende Infrastruktur zu verstehen. Durch die Problematisierung entsteht ein Bezug zu ihrer unmittelbaren Lebenswelt und ein Bewusstsein für Abhängigkeiten von Netzinfrastrukturen.

\textbf{Zukunftsbedeutung:} Kenntnisse über Netzwerke sind nicht nur für Informatik-Studiengänge oder IT-Berufe relevant, sondern allgemein für eine reflektierte Teilhabe an der digitalen Gesellschaft: Wer versteht, wie Kommunikation technisch abläuft, kann auch Fragen von Datenschutz, Netzneutralität und digitaler Souveränität kompetenter einordnen.

\textbf{Exemplarische Bedeutung:} Am Beispiel einer Chat-Kommunikation werden grundlegende Konzepte der Informatik sichtbar: IP-Adressen, Ports und Client-Server-Strukturen lassen sich am „alltäglichen Problem“ einer Chat-Nachricht erschliessen und können auf andere digitale Dienste (E-Mail, Streaming, Online-Banking) übertragen werden.

Das formulierte Problem ermöglicht es den Lernenden, technische Funktionsprinzipien mit ihrer eigenen Lebenswelt zu verbinden und diese kritisch zu reflektieren. Damit entspricht es sowohl den Prinzipien problemorientierten Lernens als auch Klafkis Anspruch auf eine didaktische Analyse.\clearpage